\documentclass{article}
\usepackage[margin=1.15in]{geometry}
\usepackage{xcolor}
\usepackage{parskip}
\usepackage{fancyhdr}
\usepackage{array}
\usepackage{booktabs}
\usepackage{enumitem}
\usepackage{amsmath}
\usepackage{tabularx}

\pagestyle{fancy}
\fancyhf{}
\cfoot{\thepage}
\rhead{Lecture 1}
\lhead{MA 214: Applied Statistics}
\renewcommand{\headrulewidth}{.4mm}

\newcommand{\blankline}[1]{\underline{\hspace{#1}}}
\newcommand{\shortanswer}{\vspace{3em}}
\newcommand{\medanswer}{\vspace{4.5em}}

\begin{document}

\noindent\textbf{Name:}\blankline{2.25in}\hfill\textbf{BUID:}\blankline{1.55in}\newline

\begin{center}
 \Large Lecture 1 In-Class Worksheet: Getting Started with MA214
\end{center}

\subsection*{Scenario}
Today is the first lecture of MA214. This worksheet is meant to help you practice the course routine and review the statistical language that will be used throughout the semester.

\medskip
\noindent\textbf{Goals:} \textbf{(1)} identify where course tasks fit in the weekly workflow, \textbf{(2)} interpret simple probability statements, \textbf{(3)} distinguish parameters, statistics, and inference goals.

\subsection*{Part 1: First-Week Course Routine}
This first part is not a statistics problem. The goal is to make sure you know the basic weekly rhythm of the course and where to find course information.

\begin{enumerate}[label={\bfseries (\Alph*)},itemsep=.8em]
\item \textbf{Weekly map.} Fill in the main course activity or deadline for each day. Some days may have more than one item.

\begin{center}
\renewcommand{\arraystretch}{1.45}
\begin{tabularx}{\linewidth}{>{\bfseries}p{2.4cm}X}
\toprule
\textbf{Day} & \textbf{What usually happens?} \\
\midrule
Monday & \\[2em]
Tuesday & \\[2em]
Wednesday & \\[2em]
Thursday & \\
\bottomrule
\end{tabularx}
\end{center}

\item \textbf{Where would you check first?} Match each situation to a reasonable first place to look or person to contact.

\begin{center}
\renewcommand{\arraystretch}{1.35}
\begin{tabularx}{\linewidth}{p{8.4cm}X}
\toprule
\textbf{Situation} & \textbf{Check first} \\
\midrule
You need the lecture slides or recording. & \\[1.6em]
You want to know when the next homework or quiz is due. & \\[1.6em]
You need to submit a lab deliverable or post-lab activity. & \\[1.6em]
You have a question about a lecture concept. & \\[1.6em]
You need to miss lecture, lab, or discussion. & \\
\bottomrule
\end{tabularx}
\end{center}

\item \textbf{This week.} Write two concrete next steps you should take after today's lecture.
\begin{enumerate}[label=\arabic*.,itemsep=1.6em]
\item \blankline{5.5in}
\item \blankline{5.5in}
\end{enumerate}
\end{enumerate}
\subsection*{Part 2: Events and Interpreting a Probability}
Suppose a weather model gives
\[
P(\text{rain tomorrow}) = 0.30.
\]
\begin{enumerate}[label={\bfseries (\Alph*)},itemsep=.75em]
\item Identify the \textbf{random process}, one possible \textbf{outcome}, and the \textbf{event} whose probability is given above. \\[3em]
\item If there are 100 days with similar weather conditions, about how many of those days would you expect to have rain? \\[2.5em]
\item If tomorrow and the day after tomorrow have similar conditions, and we assume the two days are independent, compute the probability that it does \textbf{not} rain on either day.
\[
P(\text{no rain on both days}) = \blankline{2.5in}
\]
\item If it does not rain tomorrow, does that mean the model was wrong? Explain briefly. \\[3.5em]
\end{enumerate}

\subsection*{Part 3: Parameter, Statistic, or Inference?}
Suppose we take a random sample of 100 flights departing LAX.

\begin{center}
\renewcommand{\arraystretch}{1.3}
\begin{tabular}{lcc}
\toprule
\textbf{Carrier} & \textbf{Number of sampled flights} & \textbf{Mean distance} \\
\midrule
AA & 40 & 1,220 miles \\
UA & 35 & 1,380 miles \\
Other & 25 & 1,100 miles \\
All sampled flights & 100 & 1,230 miles \\
\bottomrule
\end{tabular}
\end{center}

For each statement, decide whether it is a \textbf{parameter}, \textbf{statistic}, or \textbf{inference}. Write one short reason.

\begin{center}
\renewcommand{\arraystretch}{1.5}
\begin{tabularx}{\textwidth}{p{8.1cm}p{2.6cm}X}
\toprule
\textbf{Statement} & \textbf{Type} & \textbf{Reason} \\
\midrule
The average distance of the 100 sampled flights is 1,230 miles. & & \\[2.2em]
The true average distance of all flights departing LAX is unknown. & & \\[2.2em]
We estimate that the average distance of all LAX flights is about 1,230 miles. & & \\[2.2em]
The true average distance of all UA flights departing LAX is $\mu_{UA}$. & & \\
\bottomrule
\end{tabularx}
\end{center}

\subsection*{Part 4: Goals of Statistical Inference}
Classify each question as one of the following: \textbf{description}, \textbf{estimation}, \textbf{testing}, \textbf{prediction}, or \textbf{causation}.

\begin{center}
\renewcommand{\arraystretch}{1.45}
\begin{tabularx}{\textwidth}{p{9cm}X}
\toprule
\textbf{Question} & \textbf{Goal} \\
\midrule
What patterns appear in the 100 sampled LAX flights? & \\[1.8em]
What is the true average distance of all LAX flights? & \\[1.8em]
Is the sample mean distance for UA flights unusually large compared with the other sampled flights? & \\[1.8em]
Can we predict the distance of a future LAX flight using the carrier? & \\[1.8em]
Would changing a flight's carrier cause its distance to change? & \\
\bottomrule
\end{tabularx}
\end{center}

\subsection*{Wrap-Up}
Write one question you still have about the course workflow or today's statistics review. 
\end{document}
